\begin{textsample}{POS Dim 2 – human – Score 6.00 – rj/rj\_ef\_ciencias\_humanas\_geo\_2\_p8.txt}  \label{ex:f2_pos_017}
Considerando que o conhecimento é construído pelo sujeito – no caso, o estudante - em sua relação com os outros e com o mundo, é necessário que os conhecimentos sejam refletidos, reelaborados pelo estudante para se construir em conhecimento apropriado por ele.
Assim, é preciso desenvolver competências do pensar, ensinar a pensar, desenvolver a \textbf{capacidade} de elaboração própria. [... ] importa não apenas \textbf{buscar} os \textbf{meios} pedagógico-didáticos de melhorar e potencializar a aprendizagem pelas competências do pensar, mas também de ganhar elementos conceituais para a apropriação crítica da realidade. [... ].
Pensar é mais do que explicar, e para isso as escolas e as instituições formadoras de professores precisam formar sujeitos pensantes, \textbf{capazes} de pensar epistêmico, ou seja, sujeitos que desenvolvam \textbf{capacidades} básicas de pensamento, elementos conceituais que lhes permitam, mais do que \textbf{saber} coisas, mais do que receber uma informação, colocar -se ante a realidade, apropriar -se do momento histórico para pensar historicamente essa realidade e reagir a ela. ( ZEMELMAN, 1994, apud LIBÂNEO, 1998 )

% matched lemmas: buscar, capacidade, capaz, meio, saber
\end{textsample}
